\subsection{APP--Witt intertwining, Leech neighbor, and orbifold of order \(3\)}

The twelve fibers bear the labels
\[
 (i,\mu,\varepsilon)\in
 \mathbb Z/3\mathbb Z\times\{0,1\}\times\{0,1\},
 \qquad
 \lambda(i,\mu,\varepsilon)=4i+2\mu+\varepsilon\pmod{12}.
\]
The map \(\lambda\) is bijective and, once trivializations and frames are fixed,
the intertwiners satisfy
\[
 g_{\lambda(i,\mu,\varepsilon)}T_{i\mu\varepsilon}
 =T_{i\mu\varepsilon}C^{(-1)^\mu}.
\]
Its direct sum is an \(C_3\)-equivariant isomorphism
\[
 T_\lambda:W_{\mathrm{APP}}\otimes\mathbb C
 \xrightarrow{\ \cong\ }
 \bigoplus_{b=0}^{11}(A_{2,b}\otimes\mathbb C).
\]

Let \(N=N(A_2^{12})\), let \(v=v_\alpha\) the vector radial of norm \(54\),
and definamos
\[
 N_0=\ker\bigl(\langle\,\cdot\,,v\rangle\bmod3\bigr),
 \qquad N_v=N_0+\mathbb Z\frac v3.
\]
The lattice calculation verifies that \(N_v\cong\Lambda_{24}\), that the
isometry \(g_c\) has order three and has no nonzero fixed points, and that
\[
 \frac{g_cv-v}{3},\quad\frac{g_c^{-1}v-v}{3}\in N_0.
\]
Hence \(g_cN_v=N_v\). Its standard lift to the lattice VOA has twisted
weight \(4/3\) and type zero. The classical cyclic-orbifold theorems give
\[
 V_{(3)}:=\operatorname{Orb}_{\mathbb Z/3}
       (V_{N_v},\widehat g_c)\cong V^\natural,
 \qquad T_{3B}=t_3+12.
\]
The character is recovered by the average of the nine sectors and by
the level three uniformization:
\[
 \operatorname{ch}_{V_{(3)}}
 =\frac13\sum_{i,j\in\mathbb Z/3}Z_{i,j}=J,
\]
\[
 J=t_3+12+\frac{10\,3^9}{t_3}
       +\frac{4\,3^{14}}{t_3^2}+\frac{3^{18}}{t_3^3}.
\]
In particular,
\[
 [q]J=54+10\cdot3^9=196884=54(5\cdot729+1).
\]

\subsection{Constructions of orders \(2\) and \(3\) of the vertex operator algebra \texorpdfstring{\(V^\natural\)}{V natural}}

Fix the central extension and the Leech cocycle.
\[
 1\longrightarrow\{\pm1\}\longrightarrow\widehat\Lambda_{24}
 \longrightarrow\Lambda_{24}\longrightarrow1,
 \qquad
 V_\Lambda=M(1)\otimes\mathbb C_\varepsilon[\Lambda_{24}].
\]
The negation \(\lambda\mapsto-\lambda\) lifts to \(\theta\), and
\[
 t_2(\tau)=\operatorname{Tr}_{V_\Lambda}
  \left(\theta q^{L_0-1}\right)
 =q^{-1}\prod_{n\ge1}(1+q^n)^{-24}.
\]
The FLM presentation is
\[
 V_{(2)}:=(V_\Lambda)^+\oplus(V_\Lambda^T)^+\cong V^\natural.
\]

\begin{theorem}[Typed equivalence of the presentations of orders \(3\) and \(2\)]
\label{thm:cierre-vtm-equivalencia-voa-23}
Let
\[
 \phi_3:V_{(3)}\xrightarrow{\ \cong\ }V^\natural,
 \qquad
 \phi_2:V_{(2)}\xrightarrow{\ \cong\ }V^\natural
\]
the graded isomorphisms provided by the orbifold theorems
applied to the preceding data. Then
\[
 \boxed{\Phi_{23}:=\phi_2^{-1}\circ\phi_3:
 V_{(3)}\xrightarrow{\ \cong\ }V_{(2)}}
\]
is a graded isomorphism of vertex operator algebras. Moreover,
\[
 a\longmapsto\Phi_{23}a\Phi_{23}^{-1}
\]
identifies its automorphism groups and preserves the graded traces.
\end{theorem}

\begin{proof}
The composition of VOA isomorphisms is a VOA isomorphism and preserves the
vacuum, the conformal vector, the vertex operators, and the grading by
\(L_0\). Conjugation transports automorphisms. For every automorphism
\(a\), invariance of the trace under conjugation gives equality of the
corresponding graded series.
\end{proof}

The isomorphism \(\Phi_{23}\) depends on the identifications \(\phi_2\) and
\(\phi_3\); absolute canonicity is not asserted. The order-two construction
also preserves the exclusive content
\[
 T_{2A}=t_2+24+\frac{2^{12}}{t_2},
 \qquad
 \boxed{j=\frac{(t_2+256)^3}{t_2^2}},
\]
which completes the modular uniformization of order two.

Once \(V^\natural\) has been constructed, the three publications are stated without
unduly serializing them:
\[
 \operatorname{Aut}(V^\natural)\cong\mathbb M,
 \qquad
 \operatorname{ch}_{V^\natural}=J,
 \qquad
 g\longmapsto T_g(\tau)
 =\operatorname{Tr}_{V^\natural}\!\left(gq^{L_0-1}\right).
\]
The Moonshine theorem coordinates the graded action with genus-zero functions
and their replicability identities. The Monster action resides
in \(V^\natural\); the fact that it is not a pointwise permutation of digits is a final type
control, not the thesis of the corridor.

